\chapter{D'où vient la latence}\label{ch:ll:where-latency-comes-from}

Une étude des données de messages d'une grande bourse d'actions a établi que les courses pour trader sur une cotation périmée, ou pour l'annuler,
avant qu'un autre n'y trade, surviennent environ une fois par minute dans chacune des plus grandes actions et représentent environ un
cinquième du volume de la journée, et se décident en général en cinq à dix millionièmes de seconde (Aquilina, Budish et
O'Neill, 2022). Dans ces courses, la microseconde est une unité de compétition. Ce livre porte sur le logiciel de part et d'autre
de cette frontière~: où va chaque microseconde d'un système de trading, comment la mesurer et comment la récupérer. Ce premier chapitre
fixe le vocabulaire~: le chemin qu'un événement de marché suit à travers un système de trading, pourquoi sa durée est une
distribution plutôt qu'un nombre, comment établir un budget pour cette durée et quelles stratégies paient pour quelle partie de celle-ci.

\section{Le chemin tick-to-trade}

\begin{definition}[Latence]\label{def:ll:where-latency-comes-from:latency}
La \emph{latence}\index{latence} d'un système entre deux événements est le temps écoulé entre le premier événement
(l'arrivée d'un paquet, l'appel d'une fonction) et le second (le départ d'un ordre, le retour de la fonction), mesuré pour une occurrence donnée.
C'est une variable aléatoire~: chaque message a la sienne.
\end{definition}

Le livre~7, chapitre~18, utilisait deux \omterm{def:ll:where-latency-comes-from:latency}{latences} d'une société de trading tout entière, la \omterm{def:ll:where-latency-comes-from:latency}{latence} des données de marché et
la \omterm{def:ll:where-latency-comes-from:latency}{latence} de saisie des ordres, pour modéliser ce que voit une stratégie et le moment où ses ordres agissent. Ce livre les décompose.

\begin{definition}[Latence tick-to-trade et latence wire-to-wire]\label{def:ll:where-latency-comes-from:t2t}
La \emph{latence tick-to-trade}\index{latence tick-to-trade} d'un système de trading est le temps écoulé entre l'arrivée du message
de données de marché qui déclenche une décision et le départ de l'ordre que cette décision envoie. Il s'agit d'une
\emph{latence wire-to-wire}\index{latence wire-to-wire} lorsque les deux instants sont relevés
sur le câble réseau (le dernier bit du paquet entrant, le premier bit du paquet sortant), par un dispositif placé hors du serveur~;
c'est une \omterm{def:ll:where-latency-comes-from:latency}{latence} tick-to-trade logicielle lorsqu'ils sont relevés à l'intérieur du programme
(le paquet remis à l'application, l'ordre remis à la pile réseau).
\end{definition}

La différence compte. Le chiffre wire-to-wire inclut la carte réseau, le chemin de réception du système d'exploitation et
le chemin d'émission de la carte~; le chiffre logiciel, qu'un programme peut mesurer lui-même avec l'horloge du processeur
(chapitre~5), les exclut. Un fournisseur qui annonce une \omterm{def:ll:where-latency-comes-from:latency}{latence} sans préciser entre quels deux instants
elle est mesurée n'a pas réellement annoncé de \omterm{def:ll:where-latency-comes-from:latency}{latence}.

\begin{omfigure}
\begin{tikzpicture}[font=\footnotesize,
  st/.style={draw=omDef,fill=omDef!8,rounded corners=2pt,minimum height=0.9cm,text width=1.3cm,align=center},
  hw/.style={draw=gray,fill=gray!10,rounded corners=2pt,minimum height=0.9cm,text width=1.1cm,align=center}]
\node[hw] (nic1) at (0,0) {carte réseau};
\node[hw] (os) at (1.5,0) {noyau ou contournement};
\node[st] (dec) at (3.2,0) {décodage, arbitrage};
\node[st] (book) at (4.9,0) {carnet};
\node[st] (strat) at (6.6,0) {stratégie};
\node[st] (risk) at (8.3,0) {contrôles de risque};
\node[st] (gw) at (10.0,0) {encodage, envoi};
\node[hw] (nic2) at (11.6,0) {carte réseau};
\draw[-Stealth,thick] (nic1) -- (os); \draw[-Stealth,thick] (os) -- (dec); \draw[-Stealth,thick] (dec) -- (book); \draw[-Stealth,thick] (book) -- (strat); \draw[-Stealth,thick] (strat) -- (risk); \draw[-Stealth,thick] (risk) -- (gw); \draw[-Stealth,thick] (gw) -- (nic2);
\draw[-Stealth,thick] (-1.1,0) node[left,align=right,font=\footnotesize]{paquet\\ entrant} -- (nic1);
\draw[-Stealth,thick] (nic2) -- (12.7,0) node[right,align=left,font=\footnotesize]{ordre\\ sortant};
\draw[decorate,decoration={brace,amplitude=5pt}] (2.5,0.6) -- (10.7,0.6) node[midway,above=6pt]{tick-to-trade logiciel (le chemin critique)};
\draw[decorate,decoration={brace,amplitude=5pt,mirror}] (-0.6,-0.6) -- (12.2,-0.6) node[midway,below=6pt]{wire-to-wire, horodaté par un dispositif de capture sur le câble};
\end{tikzpicture}

\omcaption{Le chemin tick-to-trade d'un serveur de trading. Les étapes logicielles (en bleu) sont le sujet de ce livre~;
la carte réseau et les chemins de réception et d'émission du noyau (en gris) sont traités du côté du programme aux chapitres~4
et~13, et du côté matériel dans le livre~14.}\label{fig:ll:where-latency-comes-from:path}
\end{omfigure}

\begin{definition}[Chemin critique]\label{def:ll:where-latency-comes-from:hot}
Le \emph{chemin critique}\index{chemin critique} (hot path) d'un système de trading est la suite de code exécutée entre un événement de marché déclencheur
et l'ordre qui en résulte~: le code dont la durée est la \omterm{def:ll:where-latency-comes-from:t2t}{latence tick-to-trade}. Tout le reste (écriture de journaux sur disque, reporting, agrégation des risques, configuration)
s'exécute hors du chemin critique, sur d'autres threads ou à d'autres moments.
\end{definition}

La \Cref{fig:ll:where-latency-comes-from:path} est la carte de la partie~IV de ce livre, où chaque étape est construite~:
le feed handler qui décode et arbitre les paquets de la bourse (chapitre~18), le constructeur de carnet (chapitre~19),
le moteur de stratégie (chapitre~20), la barrière de risque avant négociation (chapitre~22) et la passerelle d'ordres (chapitre~21).
Les parties~I à~III fournissent de quoi est fait chaque étape~: le processeur et sa mémoire (chapitres~2 à~5), les langages
(chapitres~6 à~14) et les protocoles (chapitres~15 à~17).

Deux grandeurs physiques bornent le chemin par le bas et valent la peine d'être calculées une fois.

\begin{example}[La lumière dans la fibre, les bits sur le fil]\label{ex:ll:where-latency-comes-from:physics}
Une fibre monomode standard a un indice de groupe effectif d'environ 1.462 à une longueur d'onde de 1\,550~nm (fiche
technique d'un fabricant)~: la lumière y parcourt un mètre en $1.462/c = \qty{4.88}{\nano\second}$, si bien que les 30 mètres
de câble entre un serveur et le commutateur d'une bourse, à l'intérieur d'un même centre de données, coûtent environ \qty{146}{\nano\second} dans chaque sens.
Placer une trame sur le fil prend sa longueur en bits divisée par le débit de la ligne, auquel s'ajoutent 20 octets de préambule,
de délimiteur de début et d'intervalle inter-trames~: une trame minimale de 64 octets occupe une liaison à 10~Gb/s pendant
$84 \times 8 / 10 = \qty{67.2}{\nano\second}$ et une liaison à 25~Gb/s pendant \qty{26.9}{\nano\second}. Un chemin logiciel
de quelques microsecondes est donc dominé par le calcul et par le système d'exploitation, et non par la physique,
tant qu'il n'est pas très bien optimisé.
\end{example}

\section{La latence est une distribution}

La \omterm{def:ll:where-latency-comes-from:t2t}{latence tick-to-trade} d'un message est un tirage. Sa médiane décrit le message ordinaire~;
les messages qui coûtent de l'argent sont souvent les autres.

\begin{definition}[Latence de queue, gigue]\label{def:ll:where-latency-comes-from:tail}
La \emph{latence de queue}\index{latence de queue} (tail latency) d'un système est la partie haute de sa distribution de \omterm{def:ll:where-latency-comes-from:latency}{latence}, rapportée
sous forme de centiles élevés~: le 99\up{e} (p99), le 99.9\up{e} (p99.9), le maximum sur une période. La \emph{gigue}\index{gigue} (jitter) est la variabilité
de la \omterm{def:ll:where-latency-comes-from:latency}{latence} d'une occurrence à la suivante, rapportée sous forme d'étendue (p99 moins p50, ou un écart type)~; un système à faible gigue
est un système dont la queue est proche de sa médiane.
\end{definition}

On note $X$ une \omterm{def:ll:where-latency-comes-from:latency}{latence}, $X_i$ la \omterm{def:ll:where-latency-comes-from:latency}{latence} de l'étape $i$ d'un chemin et $q_p(X)$ son quantile d'ordre $p$, de sorte que
p99 vaut $q_{0.99}(X)$. La \omterm{def:ll:where-latency-comes-from:latency}{latence} de bout en bout d'un chemin est $X = \sum_i X_i$. Les moyennes s'additionnent~: $\E[X] = \sum_i \E[X_i]$.
Les quantiles, eux, ne s'additionnent pas, et le sens de l'erreur n'est pas fixé.

\begin{proposition}[Queues d'une somme]\label{prop:ll:where-latency-comes-from:union}
Pour toutes variables aléatoires $X_1,\dots,X_n$, quelle que soit leur dépendance, et pour des niveaux $\alpha_i > 0$ avec $\sum_i \alpha_i < 1$,
\[
q_{1-\sum_i \alpha_i}\Big(\sum_i X_i\Big) \le \sum_i q_{1-\alpha_i}(X_i).
\]
En particulier, des cibles d'étape $b_i$ atteintes au niveau $1-(1-p)/n$ garantissent une cible de bout en bout $\sum_i b_i$ au
niveau $p$.
\end{proposition}

\begin{proof}
Posons $b_i = q_{1-\alpha_i}(X_i)$, de sorte que $\P(X_i > b_i) \le \alpha_i$. Si $\sum_i X_i > \sum_i b_i$, alors $X_i > b_i$ pour
au moins un $i$~; donc $\P(\sum_i X_i > \sum_i b_i) \le \sum_i \P(X_i > b_i) \le \sum_i \alpha_i$, et le quantile d'ordre
$(1-\sum_i\alpha_i)$ de la somme est au plus $\sum_i b_i$. Avec $\alpha_i = (1-p)/n$, le niveau est $p$.
\end{proof}

La proposition ne dit rien de $q_{0.99}(\sum X_i)$ face à $\sum q_{0.99}(X_i)$, et cette comparaison peut aller
dans un sens ou dans l'autre. Elle échoue dans le sens qui compte lorsque chaque étape présente des blocages rares
qu'aucun percentile des étapes ne voit.

\begin{example}[Deux étapes saines, un chemin bloqué]\label{ex:ll:where-latency-comes-from:nonsub}
Deux étapes indépendantes durent chacune \qty{1}{\micro\second}, sauf que chacune se bloque pendant \qty{100}{\micro\second} avec
une probabilité de 0.9~\%. Le p99 de chaque étape vaut \qty{1}{\micro\second}~: ses blocages sont plus rares que 1~\%. Le chemin se bloque dès que
l'une des deux étapes se bloque, avec une probabilité de $1-0.991^2 = 1.79~\%$, si bien que son p99 vaut \qty{101}{\micro\second}, soit cinquante fois la
somme des p99 des étapes. (La même défaillance de l'additivité pour la valeur en risque explique pourquoi le livre~6 préfère
l'expected shortfall pour agréger les risques.)
\end{example}

La \Cref{fig:ll:where-latency-comes-from:ccdf} montre le même effet sur un chemin synthétique à six étapes, chaque étape avec un
corps lognormal et une probabilité de 0.3~\% de blocage d'une durée moyenne de \qty{8}{\micro\second}. Les médianes des étapes
s'additionnent presque exactement jusqu'à la médiane du chemin (\qty{2.13}{\micro\second} contre \qty{2.17}{\micro\second}), et les moyennes
s'additionnent exactement (\qty{2.32}{\micro\second}). Le chemin est bloqué avec une probabilité $1-0.997^6 \approx 1.8~\%$, si bien que son p99
(\qty{6.56}{\micro\second}) se situe dans les blocages, alors que la somme des p99 des étapes (\qty{3.59}{\micro\second}) ne
les voit pas. La borne de l'union au niveau $1-0.01/6 \approx 99.83~\%$ par étape donne \qty{28.7}{\micro\second}~: sûre,
et quatre fois trop pessimiste. Seule la distribution de bout en bout, mesurée message par message, donne la réponse.

\begin{omfigure}
\begin{tikzpicture}
\begin{axis}[width=11.5cm,height=5.6cm,xmode=log,ymode=log,log origin y=infty,xmin=1.5,xmax=60,ymin=1e-4,ymax=1.2,
  xlabel={latence de bout en bout $x$ (\si{\micro\second}, échelle logarithmique)},ylabel={$\P(X > x)$ (échelle logarithmique)},
  tick label style={font=\small},label style={font=\small},grid=major,grid style={gray!25},
  xtick={2,5,10,20,50},xticklabels={2,5,10,20,50},
  legend image code/.code={\draw[#1] (0cm,0cm) -- (0.6cm,0cm);},
  legend columns=2,legend style={font=\footnotesize,at={(0.5,-0.3)},anchor=north,draw=none,fill=none,
  /tikz/every even column/.append style={column sep=8pt}},table/col sep=comma]
\addplot[omDef,very thick] table[x=x_us,y=ccdf]{figdata/low-latency/01-where-latency-comes-from/ccdf.csv};
\addplot[black!45,dashed,thin,domain=1.5:60,samples=2] {0.01};
\addplot[omThm,very thick,ycomb] table[x=sum_p99,y=top]{figdata/low-latency/01-where-latency-comes-from/marks.csv};
\addplot[omMeth,very thick,ycomb] table[x=p99,y=top]{figdata/low-latency/01-where-latency-comes-from/marks.csv};
\addplot[omProp,very thick,dashed,ycomb] table[x=sum_union,y=top]{figdata/low-latency/01-where-latency-comes-from/marks.csv};
\legend{queue de bout en bout,niveau de 1~\%,somme des p99 des étapes,p99 de bout en bout,{borne de l'union (étapes à 99.83~\%)}}
\end{axis}
\end{tikzpicture}

\omcaption{La queue d'un chemin synthétique à six étapes dont chaque étape se bloque pendant 0.3~\% du temps. La somme des
p99 des étapes ne voit pas les blocages et sous-estime le p99 du chemin~; la borne de l'union
de la \cref{prop:ll:where-latency-comes-from:union} la surestime. Données~: \texttt{ll\_budget.composition}, 400\,000
messages simulés.}\label{fig:ll:where-latency-comes-from:ccdf}
\end{omfigure}

La même arithmétique vaut pour un chemin réel, mesuré. La \Cref{fig:ll:where-latency-comes-from:measured} chronomètre les quatre
étapes d'un chemin jouet, le décodeur et le carnet du livre~1 avec une décision triviale et un encodeur d'ordres, message par message sur
l'ordinateur portable de ce livre. Le carnet, construit sur des arbres équilibrés et une table de hachage, est l'étape la plus lente à la médiane~; le p99.9
du décodeur, du carnet et de la décision vaut plusieurs fois leur p99, et le maximum vaut des centaines de fois la
médiane~: une interruption du fil par le système d'exploitation, que cet ordinateur portable, sans cœurs isolés, ne
empêche pas (chapitre~13).

\begin{omfigure}
\begin{tikzpicture}
\begin{axis}[width=11.5cm,height=5.6cm,ybar,bar width=7pt,ymode=log,log origin y=infty,ymin=10,ymax=100000,
  symbolic x coords={decode,book,decide,encode,total},xtick=data,xticklabels={décodage,carnet,décision,encodage,total},enlarge x limits=0.14,
  ylabel={nanosecondes (échelle logarithmique)},tick label style={font=\small},label style={font=\small},
  grid=major,grid style={gray!25},
  legend columns=5,legend style={font=\footnotesize,at={(0.5,-0.2)},anchor=north,draw=none,fill=none,
  /tikz/every even column/.append style={column sep=6pt}},table/col sep=comma]
\addplot[fill=omDef!70,draw=omDef] table[x=stage,y=p50]{figdata/low-latency/01-where-latency-comes-from/measured_path.csv};
\addplot[fill=omDef!35,draw=omDef] table[x=stage,y=p99]{figdata/low-latency/01-where-latency-comes-from/measured_path.csv};
\addplot[fill=omThm!45,draw=omThm] table[x=stage,y=p999]{figdata/low-latency/01-where-latency-comes-from/measured_path.csv};
\addplot[fill=gray!40,draw=gray] table[x=stage,y=max]{figdata/low-latency/01-where-latency-comes-from/measured_path.csv};
\legend{p50,p99,p99.9,maximum}
\end{axis}
\end{tikzpicture}

\omcaption{\omterm{def:ll:where-latency-comes-from:latency}{Latence} par message des étapes d'un chemin jouet de tick-to-trade (décoder un message, l'appliquer au carnet
du livre~1, décider, encoder un ordre), sur 200 rejeux de l'échantillon de 2\,020 messages du livre~1, fixé à un seul CPU.
Mesuré sur un ordinateur portable (Intel Core Ultra 7 155H) sous WSL2, sans cœurs isolés~; chaque chiffre inclut le coût propre de la mesure de temps, soit environ
\qty{15}{\nano\second}. Données~: \texttt{bench\_path.py}.}\label{fig:ll:where-latency-comes-from:measured}
\end{omfigure}

\section{Budgets de latence wire-to-wire}

\begin{definition}[Budget de latence]\label{def:ll:where-latency-comes-from:budget}
Un \emph{budget de latence}\index{budget de latence} est une cible pour la \omterm{def:ll:where-latency-comes-from:latency}{latence} de bout en bout d'un chemin aux percentiles énoncés,
réparti en cibles pour chaque étape, à l'aune desquelles chaque modification du système est mesurée.
\end{definition}

Un budget transforme «~aller plus vite~» en une liste de nombres qu'une étape respecte ou non. Il se rédige
de l'extérieur vers l'intérieur~: ce dont la stratégie a besoin, ce que coûtent la physique et le matériel, et ce qui reste au logiciel.

\begin{method}[Rédiger un budget de latence]\label{met:ll:where-latency-comes-from:budget}
\begin{enumerate}
\item Énoncer l'exigence du point de vue de la stratégie~: la \omterm{def:ll:where-latency-comes-from:latency}{latence} au-delà de laquelle son edge disparaît, et les
percentiles qui comptent (la médiane pour les courses~; la queue pour l'exposition aux cotations périmées).
\item Retrancher les coûts fixes~: longueurs de câble et de fibre, sérialisation, commutateurs, passerelle propre à la bourse.
\item Répartir le reliquat entre les étapes logicielles proportionnellement à ce qu'elles doivent faire, à un niveau par étape plus strict
que le niveau de bout en bout (la \cref{prop:ll:where-latency-comes-from:union}).
\item Instrumenter chaque frontière d'étape (chapitre~26) et vérifier chaque version par rapport au budget, étape par étape et de bout
en bout, sur les échantillons conjoints par message (chapitre~25).
\end{enumerate}
\end{method}

\begin{dated}{2026-09}{Chiffres de tick-to-trade publiés}\label{dat:ll:where-latency-comes-from:published}
À CppCon 2017, un ingénieur d'un grand teneur de marché électronique a donné «~environ \qty{2.5}{\micro\second}~» comme un très
bon minimum de \omterm{def:ll:where-latency-comes-from:t2t}{latence wire-to-wire} pour un système de trading logiciel. L'indice de référence STAC-T0 mesure le temps écoulé depuis le
dernier bit des données entrantes nécessaires à une décision jusqu'au premier bit de l'ordre sortant, avec UDP en entrée, TCP en sortie et
pratiquement aucune logique de trading~; un rapport de juin 2024 sur une carte à matériel programmable donne un minimum de
\qty{13.9}{\nano\second} pour des trames de 507~octets. Les deux chiffres encadrent le secteur~: le logiciel en microsecondes,
le matériel en dizaines de nanosecondes (livre~14, chapitres~6 et~7).
\end{dated}

\begin{example}[Un budget de trois microsecondes]\label{ex:ll:where-latency-comes-from:budget}
Un système logiciel doit répondre en \qty{3}{\micro\second} wire-to-wire à la médiane. La carte et le chemin de réception
par contournement du noyau (kernel bypass) prennent environ \qty{0.9}{\micro\second} et le chemin d'émission \qty{0.8}{\micro\second} (les
chiffres synthétiques de la \cref{fig:ll:where-latency-comes-from:ccdf}~)~; il reste \qty{1.3}{\micro\second} pour le décodage, le carnet, la stratégie,
le risque et l'encodage. Répartir \qty{60}{\nano\second}, \qty{120}{\nano\second}, \qty{200}{\nano\second},
\qty{50}{\nano\second} et le reste en réserve, cela tient~; le carnet jouet de la \cref{fig:ll:where-latency-comes-from:measured}
seul, à environ \qty{80}{\nano\second} à la médiane mais proche d'une microseconde au p99.9, ne tient pas
dans une cible sur la queue. Le chapitre~19 en construit une qui y tient.
\end{example}

\section{Quelles stratégies paient pour quelles microsecondes}

Une course entre deux sociétés qui réagissent au même événement est gagnée par celle dont la \omterm{def:ll:where-latency-comes-from:latency}{latence} sur ce message est la plus faible. Avec
notre \omterm{def:ll:where-latency-comes-from:latency}{latence} $X_A$ et celle du concurrent $X_B$, indépendantes,
\[
\P(\text{gain}) = \P(X_A < X_B),
\]
qui dépend entièrement des deux distributions, et non d'un seul nombre. La médiane décide de la plupart des courses~; la queue
décide des autres.

\begin{example}[Médiane contre queue dans une course]\label{ex:ll:where-latency-comes-from:race}
Notre \omterm{def:ll:where-latency-comes-from:latency}{latence} a un corps lognormal de médiane \qty{2.0}{\micro\second} et une probabilité de 2~\% d'un blocage de
\qty{30}{\micro\second} en moyenne~; la médiane du concurrent est de \qty{2.2}{\micro\second}, avec des blocages de même
durée dans 0.5~\% des cas. Simulées sur un million de courses, nous gagnons 66.1~\%. Diviser notre médiane par deux nous fait gagner 98.0~\%. Diviser par deux notre
taux de blocage, qui ramène notre p99 de \qty{23.6}{\micro\second} à \qty{3.5}{\micro\second}, nous fait gagner 66.7~\%~: une amélioration importante
d'un centile que les courses voient à peine (\cref{fig:ll:where-latency-comes-from:race}).
\end{example}

\begin{omfigure}
\begin{tikzpicture}
\begin{axis}[width=11cm,height=5.4cm,xmin=1.6,xmax=2.8,ymin=0,ymax=1,ytick={0.2,0.4,0.6,0.8,1},
  xlabel={latence médiane de notre société (\si{\micro\second})~; celle du concurrent est de 2.2},ylabel={$\P(\text{gain})$},
  tick label style={font=\small},label style={font=\small},grid=major,grid style={gray!25},
  legend columns=2,legend style={font=\footnotesize,at={(0.5,-0.3)},anchor=north,draw=none,fill=none,
  /tikz/every even column/.append style={column sep=8pt}},table/col sep=comma]
\addplot[omDef,very thick] table[x=median_us,y=with_stalls]{figdata/low-latency/01-where-latency-comes-from/race.csv};
\addplot[omThm,thick,dashed] table[x=median_us,y=no_stalls]{figdata/low-latency/01-where-latency-comes-from/race.csv};
\legend{avec nos blocages de 2~\%,sans blocages}
\end{axis}
\end{tikzpicture}

\omcaption{Probabilité de gagner une course contre un concurrent de médiane \qty{2.2}{\micro\second} et 0.5~\% de blocages,
à mesure que notre médiane varie. Supprimer tous nos blocages déplace la courbe d'environ un point~; déplacer notre médiane d'un dixième de
microseconde la déplace d'environ huit. Données~: \texttt{ll\_budget.race\_curve}, 400\,000 courses simulées par
point.}\label{fig:ll:where-latency-comes-from:race}
\end{omfigure}

La queue importe là où la question n'est pas de savoir qui arrive le premier, mais combien de temps une société reste exposée. Les cotations d'un teneur de marché
restent dans le carnet pendant que sa vision du marché est périmée~; pendant la durée d'un blocage, elles peuvent être exécutées
contre elles par quiconque a vu le nouveau prix. Le coût espéré d'un blocage est le taux d'arrivées informées multiplié
par la durée du blocage et par la perte par exécution~: c'est la sélection adverse (livre~1, chapitre~1), amplifiée par la queue. Une
stratégie d'arbitrage de \omterm{def:ll:where-latency-comes-from:latency}{latence} (livre~11, chapitre~9) paie pour sa médiane~; un teneur de marché paie pour les deux, et
pour sa queue les pires jours, lorsque blocages et nouvelles surviennent ensemble. Une société qui exécute de gros ordres sur plusieurs heures
ne paie ni l'un ni l'autre à l'échelle de la microseconde~: son \omterm{def:ll:where-latency-comes-from:budget}{budget de latence} s'exprime en millisecondes, et le temps de ses ingénieurs est
mieux employé ailleurs.

\section{Tutoriel~: mesurer un chemin et composer son budget}\label{tut:ll:where-latency-comes-from:tut}

\begin{tutorial}
\textbf{Objectif.} Mesurer les étapes d'un chemin tick-to-trade jouet, message par message, puis composer et vérifier un budget avec
\texttt{firm.latbudget}. \textbf{État final~:} \cref{fig:ll:where-latency-comes-from:measured,fig:ll:where-latency-comes-from:ccdf}
et un rapport de budget.
\begin{tutsteps}
\item \textbf{Mesurer chaque étape.} Le banc d'essai rejoue l'échantillon du livre~1 et lit le compteur d'horodatage (TSC) du processeur
entre les étapes (le chapitre~2 explique le compteur et le chapitre~5 sa calibration)~; la première passe réchauffe les caches et est
écartée.
\omcode{code/low-latency/01-where-latency-comes-from/cpp/ll_path_bench.cpp}{28}{48}{Chronométrage de quatre étapes par message avec le compteur d'horodatage.}\label{lst:ll:where-latency-comes-from:bench}
\item \textbf{Exécution épinglée~:} \texttt{python bench\_path.py} compile le banc d'essai, l'exécute sur un seul processeur avec
\texttt{taskset}, et écrit \texttt{measured\_path.csv} avec un fichier \texttt{.meta} qui consigne la machine, le
compilateur, les options et la charge.
\item \textbf{Vérifier un budget.} Les cibles par étape et les cibles de bout en bout sont des données~; la vérification compose les étapes conjointement,
message par message.
\omcode{code/firm/latbudget/firm_latbudget.py}{44}{54}{Un budget vérifié sur des échantillons alignés par étape.}\label{lst:ll:where-latency-comes-from:check}
\item \textbf{Composer des étapes synthétiques} avec \texttt{ll\_budget.composition()} et comparer la somme des centiles par étape
au centile de bout en bout.
\end{tutsteps}
\textbf{Pour aller plus loin.} Exécuter le banc d'essai sans \texttt{taskset} et à côté d'une compilation de ce livre, et
comparer le p99.9~; remplacer les blocages indépendants de \texttt{ll\_budget} par des blocages qui touchent toutes les étapes d'un message
en même temps, et observer quels centiles se déplacent.
\end{tutorial}

\section{Construction~: le budget de latence}\label{bld:ll:where-latency-comes-from:latbudget}

\begin{build}
\bfield{Rôle} Le contrat auquel chaque étape du \omterm{def:ll:where-latency-comes-from:hot}{chemin critique} (hot path) de la société est mesurée, de ce chapitre au
chemin assemblé du chapitre~26.
\bfield{Interface} \texttt{Stage(name, targets)}, \texttt{Budget(stages, end\_to\_end)},
\texttt{Budget.check(samples)} renvoyant une ligne par étape et par niveau (cible, mesuré, ok),
\texttt{compose\_joint}, \texttt{compose\_independent(samples, n, seed)}, \texttt{stage\_level(p, n)},
\texttt{allocate(total, weights)}, \texttt{quantile(x, p)} (rang le plus proche).
\bfield{Règles} Les échantillons sont en nanosecondes, un tableau par étape, alignés message par message lorsqu'ils proviennent d'une même exécution~;
les chiffres de bout en bout sont composés conjointement à partir d'échantillons alignés, jamais en additionnant des centiles~; une cible par étape au niveau
$p$ pour un niveau de bout en bout $p$ est signalée comme non fiable.
\bfield{Tests d'acceptation} \texttt{code/firm/latbudget/tests/}~: la borne de l'union sur les paires indépendantes, comonotones et
contramonotones~; le p99 non sous-additif de l'\cref{ex:ll:where-latency-comes-from:nonsub}~; un budget avec une
étape en échec.
\bfield{Pour aller plus loin} Intervalles de confiance sur chaque centile mesuré (bootstrap, livre~4, chapitre~13)~; un
budget exprimé en cycles autant qu'en nanosecondes.
\end{build}

\begin{omsources}
\item M.~Aquilina, E.~Budish et P.~O'Neill, ``Quantifying the high-frequency trading `arms race''', \emph{Quarterly
Journal of Economics} 137(1), 2022, 493--564.
\item C.~Cook, ``When a microsecond is an eternity: high performance trading systems in C++'', CppCon 2017.
\item STAC, STAC-T0 benchmark reports (tick-to-trade network I/O), 2020--2024.
\item Corning, SMF-28 ULL optical fibre, product information PI1470 (2020).
\end{omsources}

\section{Exercices}

\begin{exercise}[$\star$]\label{exo:ll:where-latency-comes-from:1}
Calculer le temps de propagation d'un serveur situé à 45~m de fibre du commutateur de la bourse, dans chaque sens, pour un indice de groupe de
1.462~?
\end{exercise}

\begin{exercise}[$\star$]\label{exo:ll:where-latency-comes-from:2}
Combien de temps une trame de 128~octets occupe-t-elle une liaison de 10~Gb/s, en comptant 20~octets de préambule, de délimiteur et d'intervalle~? Et sur une
liaison de 25~Gb/s~?
\end{exercise}

\begin{exercise}[$\star$]\label{exo:ll:where-latency-comes-from:3}
Quatre étapes ont des \omterm{def:ll:where-latency-comes-from:latency}{latences} moyennes de 300, 80, 150 et \qty{500}{\nano\second}. Quelle est la \omterm{def:ll:where-latency-comes-from:latency}{latence} moyenne de bout en bout, et
que peut-on dire de sa médiane~?
\end{exercise}

\begin{exercise}[$\star\star$]\label{exo:ll:where-latency-comes-from:4}
Un chemin comporte cinq étapes indépendantes, chacune se bloquant avec une probabilité de 0.4~\% par message. Quelle est la probabilité qu'une
message rencontre au moins un blocage~? Le p99 du chemin se situe-t-il dans les blocages~?
\end{exercise}

\begin{exercise}[$\star\star$]\label{exo:ll:where-latency-comes-from:5}
Une cible de p99.9 de bout en bout de \qty{10}{\micro\second} doit être garantie par des cibles sur quatre étapes. À quel
centile chaque étape doit-elle respecter sa cible, et comment les cibles doivent-elles se combiner~?
\end{exercise}

\begin{exercise}[$\star\star$]\label{exo:ll:where-latency-comes-from:6}
Lire la \cref{fig:ll:where-latency-comes-from:race}. De combien notre médiane doit-elle baisser pour gagner autant que ce que rapporte la suppression de tous nos blocages
à une médiane de \qty{2.0}{\micro\second}~?
\end{exercise}

\begin{exercise}[$\star\star\star$]\label{exo:ll:where-latency-comes-from:7}
\emph{Programmation.} Dans \texttt{ll\_budget}, rendre les blocages communs~: lorsque le chemin se bloque (probabilité de 1.8~\% par message),
placer le blocage sur une étape tirée au hasard, et aucun sur les autres. Lesquelles des figures de
la \cref{fig:ll:where-latency-comes-from:ccdf} changent~: les p99 par étape, le p99 de bout en bout~?
\end{exercise}

\begin{exercise}[$\star\star\star$]\label{exo:ll:where-latency-comes-from:8}
\emph{Trouver la faille.} «~Le p99 du feed handler est de \qty{400}{\nano\second}, celui de notre stratégie de \qty{600}{\nano\second}
et celui de notre passerelle de \qty{300}{\nano\second}, mesurés séparément dans trois benchmarks~: notre p99 tick-to-trade est de
\qty{1.3}{\micro\second}.~»
\end{exercise}

\section{Problème~: la microseconde qui a fait perdre la course}

\begin{problem}[{Problème du week-end --- deux sociétés, une cotation périmée}]\label{pb:ll:where-latency-comes-from:1}
Notre société et un concurrent se disputent les mêmes cotations périmées. Notre \omterm{def:ll:where-latency-comes-from:latency}{latence} a un corps lognormal (médiane
\qty{2.0}{\micro\second}, dispersion 0.15) et une probabilité de 2~\% par message d'un blocage de durée exponentielle de moyenne
\qty{30}{\micro\second}~; la médiane du concurrent est de \qty{2.2}{\micro\second} et il se bloque 0.5~\% du temps.

\medskip\noindent\textbf{Partie I --- La course.}
\begin{enumerate}
\item Écrire la probabilité de gagner une course en fonction des deux lois de \omterm{def:ll:where-latency-comes-from:latency}{latence}. Que suppose-t-elle sur
leur dépendance~?
\item Sans blocages, les deux corps étant lognormaux de même dispersion, quelle est $\P(\text{gain})$ sous forme fermée~?
L'évaluer.
\item Quelle part des courses gagnons-nous avec les blocages (la simulation)~?
\item Pourquoi la réponse est-elle inférieure à la réponse sans blocage d'une valeur moindre que notre taux de blocage~?
\end{enumerate}

\medskip\noindent\textbf{Partie II --- Médiane ou queue.}
\begin{enumerate}[resume]
\item Quels sont nos p50 et p99 avec les blocages~?
\item Que rapporte la division de notre médiane par deux~? Que rapporte la division de notre taux de blocage par deux, et que fait-elle à notre p99~?
\item Supprimer tous les blocages rapporte combien~?
\item Pour quoi la société devrait-elle payer, si elle ne fait que de la course~?
\end{enumerate}

\medskip\noindent\textbf{Partie III --- Exposition.}
\begin{enumerate}[resume]
\item Nos cotations sont exposées pendant que nous sommes bloqués. Avec des blocages sur 2~\% des messages et une moyenne de
\qty{30}{\micro\second}, quelle fraction du temps notre vision du marché est-elle périmée si un message arrive toutes les
\qty{10}{\micro\second} en moyenne et que les blocages ne se chevauchent pas~?
\item Si des preneurs informés arrivent à raison de 20 par seconde et que chaque exécution sur une cotation périmée coûte un demi-tick de 0.01~\$ sur 100
actions, que coûte la péremption sur une journée de 6.5 heures~?
\item Comment la division du taux de blocage modifie-t-elle ce coût~?
\item De quelle partie de la distribution dépend l'exposition, et de quelle partie dépend la course~?
\end{enumerate}

\medskip\noindent\textbf{Partie IV --- Le budget.}
\begin{enumerate}[resume]
\item La société veut $\P(\text{gain}) \ge 0.9$ face à ce concurrent. Lire sur la
\cref{fig:ll:where-latency-comes-from:race} la médiane qu'il lui faut.
\item Écrire cela comme un budget réparti sur cinq étapes logicielles à parts égales de ce qui reste après
\qty{1.0}{\micro\second} de réseau et de noyau. Que reçoit chaque étape~?
\item À quel niveau chaque étape doit-elle respecter sa cible pour garantir un p99 de bout en bout~?
\item Énoncer le \emph{résultat nommé}~: notre probabilité de gain, et son gain à diviser la médiane par deux comparé à celui de la division par deux du
taux de blocage.
\item Pourquoi un seul chiffre («~notre \omterm{def:ll:where-latency-comes-from:latency}{latence} est de \qty{2}{\micro\second}~») ne décrit-il ni la course ni l'exposition~?
\item Quelle mesure apprendrait à la société la distribution du concurrent~?
\item Citer une stratégie pour laquelle ni la médiane ni la queue à la microseconde ne comptent.
\item En une phrase~: que rapporte un \omterm{def:ll:where-latency-comes-from:budget}{budget de latence} à une équipe d'ingénierie~?
\end{enumerate}
\end{problem}

\section{Questions d'entretien}

\begin{interviewq}[$\star$ \iqroles{développeur}]\label{iq:ll:where-latency-comes-from:1}
Quelle est la différence entre la \omterm{def:ll:where-latency-comes-from:t2t}{latence tick-to-trade} et la \omterm{def:ll:where-latency-comes-from:t2t}{latence wire-to-wire}, et comment mesureriez-vous chacune~?
\end{interviewq}

\begin{interviewq}[$\star$ \iqroles{développeur, trader}]\label{iq:ll:where-latency-comes-from:2}
Pourquoi les sociétés de trading publient-elles des p99 et des p99.9 plutôt que la \omterm{def:ll:where-latency-comes-from:latency}{latence} moyenne~?
\end{interviewq}

\begin{interviewq}[$\star\star$ \iqroles{développeur, chercheur}]\label{iq:ll:where-latency-comes-from:3}
Trois composants ont des \omterm{def:ll:where-latency-comes-from:latency}{latences} p99 de 1, 2 et \qty{3}{\micro\second}. Que peut-on dire du p99 du
pipeline~?
\end{interviewq}

\begin{interviewq}[$\star\star$ \iqroles{développeur}]\label{iq:ll:where-latency-comes-from:4}
La lumière parcourt environ \qty{5}{\nano\second} par mètre dans la fibre. Où d'autre le temps passe-t-il entre l'arrivée d'un paquet
sur un serveur et la sortie d'un ordre de celui-ci~?
\end{interviewq}

\begin{interviewq}[$\star\star$ \iqroles{trader, chercheur}]\label{iq:ll:where-latency-comes-from:5}
Votre société remporte 60~\% des courses face à son principal concurrent. Doit-elle consacrer son prochain trimestre à la médiane ou
à la queue de sa \omterm{def:ll:where-latency-comes-from:latency}{latence}~?
\end{interviewq}

\begin{interviewq}[$\star\star\star$ \iqroles{développeur}]\label{iq:ll:where-latency-comes-from:6}
Concevez la mesure qui vous permettrait d'écrire un \omterm{def:ll:where-latency-comes-from:budget}{budget de latence} pour un système de trading existant que vous n'avez pas construit.
\end{interviewq}
